\section{Ctl-2: separating estimation from actuation}

Ctl-1 is a one-shot intervention. Ctl-2 asks what happens when the estimate is
used repeatedly. The trained model produces a clean final residual state
$h_0\in\mathbb{R}^d$. The plant readout is the model's affine target head,
\[
z(h)=b_T+\beta_T^\top h.
\]
The controller uses an affine observer,
\[
\widehat z_E(h)=b_E+\beta_E^\top h,
\qquad
\widehat e_t=y^\star-\widehat z_E(h_t),
\qquad
u_t=\operatorname{clip}(K\widehat e_t,-u_{\max},u_{\max}),
\]
and the implemented plant update is
\[
h_{t+1}=h_t+u_td_D.
\]

This is an additive final-residual loop. It does not rerun the Transformer after
each edit. The target $z(h_t)$ is the model's actual target head evaluated on
the edited residual. The collateral readout is a ridge probe fit on clean
residuals; we therefore call its change a fitted nuisance-probe displacement,
not ground truth away from the clean states.

\subsection{Why the original normalization was incomplete}

Let
\[
g_T=\beta_T^\top d_D
\qquad\text{and}\qquad
g_E=\beta_E^\top d_D
\]
be the true target gain and the observer's self-consistent gain. The original
normalization fixes $g_T=g_0$. In an unclipped affine rollout, however, the
observer error follows
\[
\widehat e_{t+1}=\widehat e_t(1-Kg_E).
\]
Thus equal target authority does not imply equal loop gain. The two gains agree
only under additional geometric conditions. We report both gains and the local
pole $1-Kg_E$ for every arm.

Before clipping, this scalar observer-error mode converges if and only if
$0<Kg_E<2$. This condition does not establish full-state stability, clipped-loop
stability, or true-target tracking. Let
\[
\beta_M=\beta_T-\beta_E,
\qquad
g_M=\beta_M^\top d_D,
\qquad
\rho=\frac{g_M}{g_E},
\]
and let $m_t=z(h_t)-\widehat z_E(h_t)$. If the observer-error mode converges,
the true target error $e_t=y^\star-z(h_t)$ approaches
\[
e_\infty=-m_0-\rho\widehat e_0.
\]
For an initially calibrated step, $|\rho|$ is the relative steady tracking
error. At $\rho=-1$, the actuator has no true-target authority; below $-1$,
it moves the target opposite to the response seen by the observer. A stable
observer pole therefore does not suffice. The actuation direction must also
have little response along what the observer omits.

When $g_E>0$, a gain-matched controller can set
\[
K'=\frac{1-p_0}{g_E}
\]
for a chosen pole $p_0$. Arms with $g_E\leq0$ are reported as
sign-incompatible with a positive-gain controller; we do not conceal the problem
by reversing control authority with a negative $K$.

\subsection{Factorial design}

The primary design crosses two estimators with two directions:
\[
\{E_{\mathrm{FO}},E_{\mathrm{lifted}}\}
\times
\{D_{\mathrm{FO}},D_{\mathrm{lifted}}\}.
\]
Oracle rows and columns serve as calibration checks. For a fixed direction, the
FO-minus-lifted contrast isolates estimation. For a fixed estimator, the
FO-minus-lifted direction contrast isolates direction geometry. The diagonal
comparison remains operationally useful, but it changes both quantities.

We run unprojected directions and a diagnostic projection into the affine span
of the four clean residual-state centroids. The projection tests whether an
off-span component explains a result. It is not a manifold-preserving
intervention: the span contains arbitrary affine combinations, and a rescaling
can move far from every observed clean state.

The four binary inputs produce four distinct trajectories per training seed.
Repeated test rows are not independent observations. We use six independently
trained models as the uncertainty unit and report their minimum and maximum,
not a nominal percentile interval.

\subsection{Estimator effect with direction held fixed}

Integrated squared target error is the primary outcome,
\[
\operatorname{ISE}=\sum_{t=0}^{T}\bigl(y^\star-z(h_t)\bigr)^2.
\]
Table~\ref{tab:ctl2-factorial} gives unprojected fixed-direction contrasts. At
every nonzero interaction strength, FO ISE exceeds lifted ISE in all six seeds
with either direction fixed. Direction geometry changes the size of the effect,
especially under the first-order estimator, but it does not remove the
estimator effect.

\begin{table}[t]
\centering
\small
\begin{tabular}{@{}lrrr@{}}
\toprule
$\gamma$ & Fixed direction & Mean FO--lifted ISE & Seed min--max \\
\midrule
0.50 & first order & 5.35 & 0.0134--24.26 \\
0.50 & lifted & 7.63 & 2.79--21.24 \\
1.00 & first order & 43.03 & 9.09--78.16 \\
1.00 & lifted & 36.05 & 14.9--65.45 \\
1.15 & first order & 60.74 & 11.7--133 \\
1.15 & lifted & 50.02 & 19--102.3 \\
1.50 & first order & 95.73 & 15.8--211 \\
1.50 & lifted & 77.09 & 25--175.9 \\
\bottomrule
\end{tabular}
\caption{Ctl-2 estimator effects with direction fixed. All listed contrasts are positive in all six training seeds. The $\gamma=0$ contrasts are numerical zero and are omitted.}
\label{tab:ctl2-factorial}
\end{table}


\begin{figure}[t]
  \centering
  \includegraphics[width=0.98\linewidth]{revision_ctl2_factorial_ise_by_gamma.png}
  \caption{Ctl-2 factorial ISE effects. Bands are the minimum and maximum over
  six training seeds. The projected panel is a support diagnostic, not an
  on-manifold result.}
  \label{fig:ctl2-factorial}
\end{figure}

The old thresholded divergence flag gives a different impression at
$\gamma=0.5$. In the unprojected diagonal comparison, its mean FO-minus-lifted
rate is $-0.151$, with seed values from $-0.456$ to $0.451$. The continuous
ISE difference is positive in all six seeds, with mean $6.71$ and range
$1.13$--$26.20$. The sign flip is a threshold artifact in a coarse flag.
We retain the flag as a secondary summary and do not describe the clipped loop
as unbounded instability.

\subsection{Gain matching and response calibration}

The no-clipping gain-matched control isolates controller gain from target
tracking. At $\gamma=0$ and a relative setpoint offset of $0.01$, the
first-order arm uses $K'\approx14.67$, realizes local pole $0.7025$, and has
zero clipping. Its target MSE nevertheless grows from $0.0001$ to $0.170$,
with ISE $1.92$. Matching the observer-error pole is not enough because it
does not correct initial bias or the relationship between observer response and
true target response.

We therefore calibrate finite responses separately from absolute levels. For
held-out probes anchored at a clean baseline $h_0$, define
\[
\delta\widehat z=\widehat z(h_0+\delta h)-\widehat z(h_0),
\qquad
\delta z=z(h_0+\delta h)-z(h_0).
\]
A through-origin scalar fit estimates the response gain, and the calibrated
observer preserves the initial estimate while scaling only its response:
\[
\widehat z_{\mathrm{cal}}(h)=\widehat z(h_0)
+a\left[\widehat z(h)-\widehat z(h_0)\right].
\]
In this affine fixture, response calibration aligns $g_E$ with $g_T$ and
removes the dynamic sign/gain error. At $\gamma=1.15$, seed 0, the calibrated
unprojected FO/FO arm ends at target MSE $0.069$ from initial MSE $0.720$.
The remaining error follows the initial estimation bias.

The calibration can be ill-conditioned. At $\gamma=0$, unprojected
first-order response scales range from about 24 to 73 across seeds. At other
interaction strengths, some estimator--direction pairs are nearly
response-orthogonal and require gains above $10^3$ or a sign reversal. A
statement that one scalar identifies the response is incomplete without the
conditioning of that scalar.

\begin{figure}[t]
  \centering
  \includegraphics[width=0.98\linewidth]{revision_ctl2_response_calibration_by_gamma.png}
  \caption{Response calibration. Scalar calibration repairs the local response
  mismatch in the affine fixture, while initial estimation bias remains.}
  \label{fig:ctl2-calibration}
\end{figure}

\subsection{Affine-span projection is not a global repair}

At $\gamma=0$, affine-span projection repairs target tracking in seed 0, but
it multiplies direction norm by about $7.5$. Cumulative fitted nuisance-probe
displacement rises from about $0.87$ to $13.36$, and the trajectory remains
farther than one clean-state RMS radius from the nearest observed residual on
average. The result diagnoses an off-span target-head component. It does not
establish a safe or on-manifold actuator.

The Ctl-2 result is therefore precise. The original diagonal comparison mixes
estimation and direction. The factorial shows that estimation misspecification
drives the main error at nonzero interaction, while direction geometry changes
its severity. In this affine plant, a held-out scalar response gain repairs the
dynamic mismatch but not initial bias, and the required gain may be poorly
conditioned.
